まず，KAKENデータベース\cite{kaken1}より，国立大学による報告書を収集した．本研究ではこのうちランダムに選択した4大学の報告書からさらにランダムに各大学につき50件を抽出し，対象とした．地名の抽出対象は研究課題名および報告書本文とした．
次にmecab\cite{mecab}により固有表現抽出（形態素解析）を行い，"固有名詞"および"地域"のラベルがあるものを抽出した．これらに対し，さらにmecabの判定が妥当かを複数の判定者に判定させた．判定者は3種のAIと3名の人とした．
判定者へは，AI，人ともに表\ref{tab:prompt}（視認性のため文には改行・空白を挿入している）の指示を与えた．データは4要素からなるJSON形式で与え，第2要素を抽出された地名，第4要素を研究課題名および本文（すなわち固有表現抽出の対象）とした．それ以外の要素は管理目的で挿入したものであり無視するよう促した．なお，AIには，Gemma3（4b）\cite{gemma3}，Llama3（8b）\cite{llama3}，Qwen3（32b）\cite{qwen3}を用いた．期待する出力形式は，入力されたJSONの第2要素に対する修正がなされたJSONである．AIの判定実行はローカルAI環境（Ollama\cite{ollama}）により1報告書ごとに行った．判定はAIに関しては一回のみ実行し，人に対しては期限を設け，複数回の提出を認めた．
最後に，判定結果を集計し，分析した．
